% ===========================================================================
% Section 7: Mechanisms and Heterogeneity
%
% Robot-mode (Shapiro Step 3): state facts, do not argue.
%
% Anchors: Table 14 (tab:mechanisms); heterogeneity results from
% code/analysis/diagnostic_heterogeneity.R (diagnostic, not a numbered
% table -- see comment before 7.2).
%
% PROVISIONAL FRAMING (decisions log F2): mechanism-first. Lead with the
% decomposition (half the MA effect runs through local infrastructure),
% give the robustness reading as the complementary interpretation in the
% following paragraph. Reversal = reorder two paragraphs.
%
% Numbers use scalars.tex macros (AutoFill pass), sourced from
% results/tables/table_14_mechanisms.csv and diagnostic_heterogeneity.csv.
%
% CITATION KEYS: none new beyond Sections 1-6.
% ===========================================================================

\section{Mechanisms and Heterogeneity}
\label{sec:mechanisms}

Market access summarizes a district's transport position relative to
every other district. A change in a district's market access can
therefore reflect infrastructure built or removed within the district
itself, or changes elsewhere in the network that alter the district's
effective distance to population. This section separates the two,
and then asks whether the population response varies with baseline
district characteristics.

\subsection{Local Infrastructure}
\label{sec:local_infra}

We measure within-district infrastructure change with four variables,
denoted $Z_i$: the change in rail kilometers between 1960 and 1986;
the change in paved and gravel road kilometers between 1954 and 1986;
an indicator for districts that lost their entire rail network by
1986 (14 districts); and an indicator for districts that gained their
first paved or gravel road (82 districts). The kilometer measures are computed
from the same georeferenced networks as the market access measures.

Table~\ref{tab:mechanisms} adds the $Z_i$ progressively to the
headline specification of Table~\ref{tab:population_iv}, column~(4).
The market-access regressor remains instrumented by both instruments;
the $Z_i$ enter as controls and are not instrumented. Column~(1)
reproduces the baseline estimate of \mainPopCoefIVBoth{}
(\mainPopSEIVBoth{}). Adding the two kilometer measures in column~(2)
lowers the market-access coefficient to \mechKmCoef{} (\mechKmSE{});
adding the two binary margins alone in column~(3) leaves it at
\mechBinCoef{} (\mechBinSE{}); and including all four $Z_i$ jointly
in column~(4) lowers it to \mechAllZCoef{} (\mechAllZSE{}). The
first-stage $F$ is higher in every augmented column (up to
\mechMaxAugF{}) than at baseline (\mechBaselineF{}), so the attenuation
is not a weak-instrument artifact. Columns~(2)--(4) lose three
districts with missing kilometer measures (308 versus 311
observations); re-estimated on the common 308-district sample, the
baseline coefficient is \mechBaselineCommonCoef{}
(\mechBaselineCommonSE{}), so the sample change does not generate the
attenuation and, if anything, works against it.

Read as a decomposition, column~(4) indicates that roughly half of
the headline market-access coefficient operates through
infrastructure located in the district itself: conditioning on own
rail and road change and the two binary margins removes about half of
the point estimate. The change in own rail kilometers enters
positively and precisely, at \mechRailKmCoef{} (\mechRailKmSE{}) in
column~(4): districts that lost more rail kilometers grew more
slowly, conditional on market access. The change in own road
kilometers enters with a coefficient indistinguishable from zero.

The same numbers support a complementary reading. Even conditional on
every measured form of own-district infrastructure change, half the
market-access coefficient remains. Market access is therefore not
merely a repackaging of local infrastructure counts; it carries
information about regional connectivity beyond the district's own
kilometers. The two readings describe the same estimate and are not
alternatives between which the data can select.

Two caveats apply. First, because the $Z_i$ are not instrumented,
column~(4) is a descriptive decomposition, not a causal mediation
analysis: own-infrastructure change is itself an outcome of the same
reform. Second, the indicator for losing the entire rail network
enters positively, at \mechLostRailCoef{} (\mechLostRailSE{}), in
column~(4). With
only 14 districts in that group, we do not interpret this
coefficient beyond noting it as a candidate for future work with
richer within-district data.

\subsection{Heterogeneity}
\label{sec:heterogeneity}

% NOTE: results below come from code/analysis/diagnostic_heterogeneity.R
% (results/tables/diagnostic_heterogeneity.{txt,csv}), run as a
% diagnostic rather than a numbered table while Block 2 framing is
% provisional. Lift into a numbered table if the coauthor meeting
% confirms the section structure.

If market access matters more where initial conditions leave more
room for reallocation, its coefficient should vary with baseline
characteristics. We examine three, one at a time, each standardized
within the estimation sample: initial log population in 1960, the
rural population share in 1960, and distance to Buenos Aires. Under
the reallocation hypothesis the interaction should be negative for
initial population and positive for the rural share and for distance
to Buenos Aires. Each specification adds the interaction of the
characteristic with $\Delta \ln \mathrm{MA}^{\mathrm{full}}_i$ to the
main equation, instrumenting the interaction with the interactions of
both instruments with the characteristic.\footnote{Estimates in this
subsection are produced by the replication archive's heterogeneity
program and are reported in its output files; a tabulated version is
deferred until the section structure is final.}

The OLS interaction coefficients carry the predicted sign in all
three cases: $\heteroPopIntOLS{}$ (\heteroPopIntSEOLS{}) for initial
population, \heteroRurIntOLS{} (\heteroRurIntSEOLS{}) for the rural
share, and an imprecise \heteroDistIntOLS{} (\heteroDistIntSEOLS{})
for distance to Buenos Aires. Under instrumental variables the signs are
unchanged but the estimates attenuate, and the first stage for the
market-access level weakens to an $F$ near 7 once the interaction
enters, because the instrument set must now separate two endogenous
regressors. We therefore read these results as sign patterns
consistent with the market-access channel mattering most in
agricultural, less urbanized, and remote districts, and not as
magnitudes. A design with stronger identification of the interactions
is left for future work.
